% GENERATED FILE. Edit canonical lessons, then run npm run generate:books.
% canonical-source-hash: fnv1a64:4a0c647b4c7c6dba
% canonical-lessons: PT-C489-insonia, PT-C489-ressonar, PT-C489-bocejar, PT-C489-arrepio, PT-C489-caibra

\chapter{Noites mal dormidas}
\label{ch:pt-noites-mal-dormidas}
\hlchaptermodality{489}

\begin{chapteropening}
\textbf{By the end of this chapter:} \emph{I can describe bad nights: snoring, yawning and cramps.}

Complete the last lesson of chapter 489: I can describe bad nights: snoring, yawning and cramps.
\end{chapteropening}

\section[a insónia]{a insónia — insomnia, sleeplessness}

\label{lesson:PT-C489-insonia}



\begin{quote}
Before the new one: say the Portuguese for to recover, then the Portuguese for recovery.
\end{quote}

\subsection*{You'll want to know: a insónia}
\textbf{a insónia} — \textquotedblleft{}insomnia, sleeplessness\textquotedblright{}.

\begin{quote}
O café à noite dá-me insónias.
\end{quote}

\begin{quote}
\emph{Coffee at night keeps me awake.}
\end{quote}

In Brazil: \emph{insônia}. Often plural: \emph{ter insónias}.

\begin{cousinweb}
\emph{insónia} is a learned borrowing from Latin \emph{īnsomnia} \textquotedblleft{}sleeplessness\textquotedblright{} (Lewis \& Short), built on \emph{somnus} \textquotedblleft{}sleep\textquotedblright{}.
\end{cousinweb}

\begin{grammarlens}[title={What you've built}]
One more word for telling, explaining and describing at B1.
\end{grammarlens}

\subsection*{Your turn}
\begin{itemize}
\raggedright
  \item \emph{Say it:} \emph{a insónia}
  \item \emph{Say it:} \emph{a insónia}, once more
  \item \emph{Say it:} a sentence of your own with \emph{a insónia}
  \item \emph{Recall:} say the Portuguese for to recover, then the Portuguese for recovery, then say \emph{a insónia} again
\end{itemize}

\begin{tcolorbox}[breakable,colback=teal!4,colframe=teal!35!black,title={Before you move on}]
What does a insónia mean? (\textquotedblleft{}Insomnia, sleeplessness\textquotedblright{}.) Say it once more.
\end{tcolorbox}
\section[ressonar]{ressonar — to snore}

\label{lesson:PT-C489-ressonar}



\begin{quote}
Before the new one: say the Portuguese for recovery, then the Portuguese for insomnia.
\end{quote}

\subsection*{You'll want to know: ressonar}
\textbf{ressonar} — \textquotedblleft{}to snore\textquotedblright{}.

\begin{quote}
O meu marido ressona tanto que não consigo dormir.
\end{quote}

\begin{quote}
\emph{My husband snores so much that I cannot sleep.}
\end{quote}

In Brazil: \emph{roncar}.

\begin{cousinweb}
\emph{ressonar} comes from Latin \emph{resonāre} \textquotedblleft{}to resound, to echo\textquotedblright{} (EtymDB), built on \emph{sonāre} \textquotedblleft{}to sound\textquotedblright{}.
\end{cousinweb}

\begin{grammarlens}[title={What you've built}]
One more word for telling, explaining and describing at B1.
\end{grammarlens}

\subsection*{Your turn}
\begin{itemize}
\raggedright
  \item \emph{Say it:} \emph{ressonar}
  \item \emph{Say it:} \emph{ressonar}, once more
  \item \emph{Say it:} a sentence of your own with \emph{ressonar}
  \item \emph{Recall:} say the Portuguese for recovery, then the Portuguese for insomnia, then say \emph{ressonar} again
\end{itemize}

\begin{tcolorbox}[breakable,colback=teal!4,colframe=teal!35!black,title={Before you move on}]
What does ressonar mean? (\textquotedblleft{}To snore\textquotedblright{}.) Say it once more.
\end{tcolorbox}
\section[bocejar]{bocejar — to yawn}

\label{lesson:PT-C489-bocejar}



\begin{quote}
Before the new one: say the Portuguese for insomnia, then the Portuguese for to snore.
\end{quote}

\subsection*{You'll want to know: bocejar}
\textbf{bocejar} — \textquotedblleft{}to yawn\textquotedblright{}.

\begin{quote}
Os alunos bocejavam na aula de sexta à tarde.
\end{quote}

\begin{quote}
\emph{The pupils were yawning in the Friday afternoon class.}
\end{quote}

A yawn is \emph{um bocejo}.

\begin{cousinweb}
\emph{bocejar} is usually explained as built on \emph{boca} \textquotedblleft{}mouth\textquotedblright{}, which is inherited from Latin \emph{bucca} \textquotedblleft{}cheek, mouth\textquotedblright{} (Lewis \& Short).
\end{cousinweb}

\begin{grammarlens}[title={What you've built}]
One more word for telling, explaining and describing at B1.
\end{grammarlens}

\subsection*{Your turn}
\begin{itemize}
\raggedright
  \item \emph{Say it:} \emph{bocejar}
  \item \emph{Say it:} \emph{bocejar}, once more
  \item \emph{Say it:} a sentence of your own with \emph{bocejar}
  \item \emph{Recall:} say the Portuguese for insomnia, then the Portuguese for to snore, then say \emph{bocejar} again
\end{itemize}

\begin{tcolorbox}[breakable,colback=teal!4,colframe=teal!35!black,title={Before you move on}]
What does bocejar mean? (\textquotedblleft{}To yawn\textquotedblright{}.) Say it once more.
\end{tcolorbox}
\section[o arrepio]{o arrepio — shiver; goosebumps}

\label{lesson:PT-C489-arrepio}



\begin{quote}
Before the new one: say the Portuguese for to snore, then the Portuguese for to yawn.
\end{quote}

\subsection*{You'll want to know: o arrepio}
\textbf{o arrepio} — \textquotedblleft{}shiver; goosebumps\textquotedblright{}.

\begin{quote}
Aquela história deu-me arrepios.
\end{quote}

\begin{quote}
\emph{That story gave me the shivers.}
\end{quote}

\emph{Dar arrepios} — to give the creeps. The verb is \emph{arrepiar}.

\begin{cousinweb}
\emph{arrepio} comes from \emph{arrepiar} \textquotedblleft{}to make one's hair stand on end\textquotedblright{}, which is usually traced to Latin \emph{horripilāre} \textquotedblleft{}to bristle with hairs\textquotedblright{} (EtymDB; Lewis \& Short).
\end{cousinweb}

\begin{grammarlens}[title={What you've built}]
One more word for telling, explaining and describing at B1.
\end{grammarlens}

\subsection*{Your turn}
\begin{itemize}
\raggedright
  \item \emph{Say it:} \emph{o arrepio}
  \item \emph{Say it:} \emph{o arrepio}, once more
  \item \emph{Say it:} a sentence of your own with \emph{o arrepio}
  \item \emph{Recall:} say the Portuguese for to snore, then the Portuguese for to yawn, then say \emph{o arrepio} again
\end{itemize}

\begin{tcolorbox}[breakable,colback=teal!4,colframe=teal!35!black,title={Before you move on}]
What does o arrepio mean? (\textquotedblleft{}Shiver; goosebumps\textquotedblright{}.) Say it once more.
\end{tcolorbox}
\section[a cãibra]{a cãibra — cramp}

\label{lesson:PT-C489-caibra}



\begin{quote}
Before the new one: say the Portuguese for to yawn, then the Portuguese for shiver.
\end{quote}

\subsection*{You'll want to know: a cãibra}
\textbf{a cãibra} — \textquotedblleft{}cramp\textquotedblright{}.

\begin{quote}
A meio da noite, acordei com uma cãibra na perna.
\end{quote}

\begin{quote}
\emph{In the middle of the night, I woke up with a cramp in my leg.}
\end{quote}

In Brazil also \emph{câimbra}.

\begin{cousinweb}
\emph{cãibra} is usually traced to a Germanic word meaning \textquotedblleft{}bent, crooked\textquotedblright{}, the same source as English \emph{cramp}; the path into Portuguese is not settled.
\end{cousinweb}

\begin{grammarlens}[title={What you've built}]
One more word for telling, explaining and describing at B1.
\end{grammarlens}

\subsection*{Your turn}
\begin{itemize}
\raggedright
  \item \emph{Say it:} \emph{a cãibra}
  \item \emph{Say it:} \emph{a cãibra}, once more
  \item \emph{Say it:} a sentence of your own with \emph{a cãibra}
  \item \emph{Recall:} say the Portuguese for to yawn, then the Portuguese for shiver, then say \emph{a cãibra} again
\end{itemize}

\begin{tcolorbox}[breakable,colback=teal!4,colframe=teal!35!black,title={Before you move on}]
What does a cãibra mean? (\textquotedblleft{}Cramp\textquotedblright{}.) Say it once more.
\end{tcolorbox}
